% GENERATED FILE. Edit canonical lessons, then run npm run generate:books.
% canonical-source-hash: fnv1a64:29642ce49ae05133
% canonical-lessons: GE-C173-ziemlich, GE-C173-besonders, GE-C173-ungefaehr, GE-C173-zusammen, GE-C173-allein, GE-R173-first-pass-words-from-chapters-168-170, GE-R173-second-pass-words-from-chapters-171-173

\chapter{Quite, Especially, About, Together, Alone}
\label{ch:ge-quite-especially-about-together-alone}
\hlchaptermodality{173}

\begin{chapteropening}
\textbf{By the end of this chapter:} \emph{I can say quite, especially, about, together and alone.}

Complete the last lesson of chapter 173: I can say quite, especially, about, together and alone.
\end{chapteropening}

\section[ziemlich]{ziemlich — quite, fairly}

\label{lesson:GE-C173-ziemlich}



\begin{quote}
Before the new one: say the German for of course, then the German for exactly.
\end{quote}

\subsection*{You'll want to know: ziemlich}
\textbf{ziemlich} — \textquotedblleft{}quite, fairly\textquotedblright{}.

\begin{grammarlens}[title={What you've built}]
One more word for this chapter.
\end{grammarlens}

\subsection*{Your turn}
\begin{itemize}
\raggedright
  \item \emph{Say it:} \emph{ziemlich}
  \item \emph{Say it:} \emph{ziemlich}, once more
  \item \emph{Say it:} say \emph{ziemlich}
  \item \emph{Recall:} say the German for of course, then the German for exactly, then say \emph{ziemlich} again
\end{itemize}

\begin{tcolorbox}[breakable,colback=teal!4,colframe=teal!35!black,title={Before you move on}]
What does ziemlich mean? (\textquotedblleft{}Quite, fairly\textquotedblright{}.) Say it once more.
\end{tcolorbox}
\section[besonders]{besonders — especially}

\label{lesson:GE-C173-besonders}



\begin{quote}
Before the new one: say the German for exactly, then the German for quite.
\end{quote}

\subsection*{You'll want to know: besonders}
\textbf{besonders} — \textquotedblleft{}especially\textquotedblright{}.

\begin{grammarlens}[title={What you've built}]
One more word for this chapter.
\end{grammarlens}

\subsection*{Your turn}
\begin{itemize}
\raggedright
  \item \emph{Say it:} \emph{besonders}
  \item \emph{Say it:} \emph{besonders}, once more
  \item \emph{Say it:} say \emph{besonders}
  \item \emph{Recall:} say the German for exactly, then the German for quite, then say \emph{besonders} again
\end{itemize}

\begin{tcolorbox}[breakable,colback=teal!4,colframe=teal!35!black,title={Before you move on}]
What does besonders mean? (\textquotedblleft{}Especially\textquotedblright{}.) Say it once more.
\end{tcolorbox}
\section[ungefähr]{ungefähr — about, roughly}

\label{lesson:GE-C173-ungefaehr}



\begin{quote}
Before the new one: say the German for quite, then the German for especially.
\end{quote}

\subsection*{You'll want to know: ungefähr}
\textbf{ungefähr} — \textquotedblleft{}about, roughly\textquotedblright{}.

\begin{grammarlens}[title={What you've built}]
One more word for this chapter.
\end{grammarlens}

\subsection*{Your turn}
\begin{itemize}
\raggedright
  \item \emph{Say it:} \emph{ungefähr}
  \item \emph{Say it:} \emph{ungefähr}, once more
  \item \emph{Say it:} say \emph{ungefähr}
  \item \emph{Recall:} say the German for quite, then the German for especially, then say \emph{ungefähr} again
\end{itemize}

\begin{tcolorbox}[breakable,colback=teal!4,colframe=teal!35!black,title={Before you move on}]
What does ungefähr mean? (\textquotedblleft{}About, roughly\textquotedblright{}.) Say it once more.
\end{tcolorbox}
\section[zusammen]{zusammen — together}

\label{lesson:GE-C173-zusammen}



\begin{quote}
Before the new one: say the German for especially, then the German for about.
\end{quote}

\subsection*{You'll want to know: zusammen}
\textbf{zusammen} — \textquotedblleft{}together\textquotedblright{}.

\begin{grammarlens}[title={What you've built}]
One more word for this chapter.
\end{grammarlens}

\subsection*{Your turn}
\begin{itemize}
\raggedright
  \item \emph{Say it:} \emph{zusammen}
  \item \emph{Say it:} \emph{zusammen}, once more
  \item \emph{Say it:} say \emph{zusammen}
  \item \emph{Recall:} say the German for especially, then the German for about, then say \emph{zusammen} again
\end{itemize}

\begin{tcolorbox}[breakable,colback=teal!4,colframe=teal!35!black,title={Before you move on}]
What does zusammen mean? (\textquotedblleft{}Together\textquotedblright{}.) Say it once more.
\end{tcolorbox}
\section[allein]{allein — alone}

\label{lesson:GE-C173-allein}



\begin{quote}
Before the new one: say the German for about, then the German for together.
\end{quote}

\subsection*{You'll want to know: allein}
\textbf{allein} — \textquotedblleft{}alone\textquotedblright{}.

\begin{grammarlens}[title={What you've built}]
One more word for this chapter.
\end{grammarlens}

\subsection*{Your turn}
\begin{itemize}
\raggedright
  \item \emph{Say it:} \emph{allein}
  \item \emph{Say it:} \emph{allein}, once more
  \item \emph{Say it:} say \emph{allein}
  \item \emph{Recall:} say the German for about, then the German for together, then say \emph{allein} again
\end{itemize}

\begin{tcolorbox}[breakable,colback=teal!4,colframe=teal!35!black,title={Before you move on}]
What does allein mean? (\textquotedblleft{}Alone\textquotedblright{}.) Say it once more.
\end{tcolorbox}
\section[(dialogue)]{Words from chapters 168-170 — a first pass}

\label{lesson:GE-R173-first-pass-words-from-chapters-168-170}



\begin{quote}
Say the words for together and alone.
\end{quote}

\subsection*{Your turn}
Say these aloud:

\begin{itemize}
\raggedright
  \item impossible, bright, dark, full, old-fashioned
  \item well known, foreign, nervous, proud, angry
  \item rarely, mostly, right away, finally, still
\end{itemize}

\begin{tcolorbox}[breakable,colback=teal!4,colframe=teal!35!black,title={Before you move on}]
Impossible? (\textbf{unmöglich}.) Bright? (\textbf{hell}.) Dark? (\textbf{dunkel}.) Full? (\textbf{voll}.) Old-fashioned? (\textbf{altmodisch}.) Well known? (\textbf{bekannt}.) Foreign? (\textbf{fremd}.) Nervous? (\textbf{nervös}.) Proud? (\textbf{stolz}.) Angry? (\textbf{wütend}.) Rarely? (\textbf{selten}.) Mostly? (\textbf{meistens}.) Right away? (\textbf{gleich}.) Finally? (\textbf{endlich}.) Still? (\textbf{noch}.)
\end{tcolorbox}
\section[(dialogue)]{Words from chapters 171-173 — a second pass}

\label{lesson:GE-R173-second-pass-words-from-chapters-171-173}



\begin{quote}
Say the words for only and first.
\end{quote}

\subsection*{Your turn}
Say these aloud:

\begin{itemize}
\raggedright
  \item only, first, afterwards, before, later
  \item unfortunately, perhaps, probably, of course, exactly
  \item quite, especially, about, together, alone
\end{itemize}

\begin{tcolorbox}[breakable,colback=teal!4,colframe=teal!35!black,title={Before you move on}]
Only? (\textbf{erst}.) First? (\textbf{zuerst}.) Afterwards? (\textbf{danach}.) Before? (\textbf{vorher}.) Later? (\textbf{nachher}.) Unfortunately? (\textbf{leider}.) Perhaps? (\textbf{vielleicht}.) Probably? (\textbf{wahrscheinlich}.) Of course? (\textbf{natürlich}.) Exactly? (\textbf{genau}.) Quite? (\textbf{ziemlich}.) Especially? (\textbf{besonders}.) About? (\textbf{ungefähr}.) Together? (\textbf{zusammen}.) Alone? (\textbf{allein}.)
\end{tcolorbox}
